\section{Educational Plan}

\subsection{Educational Activities Undertaken}

One of my main teaching goals upon my arrival at
U.C. Santa Barbara is to expand and modernize the curriculum in 
multimedia systems. 
As a first step, I will design a multimedia system course
that surveys the related disciplines and gives students opportunities
to participate in a few small-scale projects, such as
investigating feature representations for various multimedia
data types and studying caching policies for different data objects. 
I intend to keep the contents of this new multimedia system 
course I am designing quite dynamic.
While some of the topics that I will teach are mature and 
reasonably stable (e.g., image compression and segmentation), many of the 
others are evolving fast. I will use textbooks to cover the more
mature topics and will teach the rest of the topics through research
papers. This way, I expect that
this advanced course can be up-to-date and
beneficial to both students and faculty and that it will
allow us to identify interesting
and open areas that might lead to research activities. 
This course will also provide students with valuable knowledge
and skills for the job market.

To discuss more recent research results and initiatives, I
will design and lead a graduate seminar in 
multimedia databases and systems, where faculty 
and students will discuss
recent publications and industry standards (e.g., MPEG-4, MPEG-7,
ATVEF, etc.). I view this seminar as a way to start
my teaching and research career at U.C. Santa Barbara.
It will give me a perfect opportunity to meet
colleagues and students with similar interests and to
start fruitful interactions with them.

Finally, I plan to sponsor undergraduate research projects.
Both of my research focuses,
image/video databases and multimedia
systems, are very popular subjects, and I expect that
the student population in general, and undergraduate students
in particular, will be very interested in both. I envision
these projects to have a research and an implementation
component, and I expect them to show undergraduates
how research can be intellectually exciting and at
the same time practical in solving real-world problems. Students
will be exposed to cutting edge research topics and technology,
and will contribute to a larger research and implementation
effort. Both will provide valuable experience
for their future professional careers.

\subsection{Prior Teaching Experience}

I believe that being a good teacher requires not only the desire to
pass knowledge on to others, but also the ability to
cultivate and inspire young minds to discover the unknown
and to expand human possibilities.
Teaching is a privilege that also carries
great responsibilities. At the very least, a teacher must
present the course materials clearly.  More important,
a teacher must motivate the interest
of students with different backgrounds and levels of
maturity and must inspire them to progress
beyond what they think they can do. 

To prepare myself to teach effectively, I have actively sought out
opportunities to improve my teaching skills.
As a graduate student, I was a teaching assistant for an undergraduate 
``Introduction to Programming Languages'' course and for a graduate 
``Advanced Topics in Operating Systems'' course.
I prepared and graded assignments and solutions, and assisted
in preparing and grading exams. I also led discussion
sessions, which proved a useful way to learn to address the
students' weak points directly.
                                  
In addition, I had significant teaching experience as
the instructor of a database system course in the
summer of 1998 at Stanford. 
The course was broadcast over Stanford's Instructional
Television Network. Interacting with forty students,
both local and remote,
and with both full-time and part-time students working in Silicon
Valley companies presented challenging problems and
opportunities because the students' backgrounds, motivation,
and dedication varied greatly.
Directing this class was a wonderful experience. In particular,
I enjoyed the interaction with students, which gave me
the chance to learn about their needs and to assist them on an
individual basis. In this course, I also used the first draft of a new
textbook written by professors Garcia-Molina, Ullman,
and Widom; I gathered and provided feedback to the
professors for revisions before the first edition came out.

I believe that my experience has prepared me for the teaching
requirements of an academic position. Moreover,
my desire to teach well and my patience should
help me remain a very good teacher.

